%==============================================================================
%  Bachelor's Thesis Project Proposal — ETH Zurich registration
%  Single file, self-contained: compiles with pdflatex, no external .bib.
%
%  TITLE OPTIONS (chosen one is in \title; the rest kept so the decision stays
%  open until registration).
%
%  A. Cooperative Multi-Agent Generative Jamming of UAV Networks under
%     Detection Constraints
%       -> Descriptive and safe: names the method (MARL + generative), the
%          target (UAV networks) and the constraint (detection). Recommended
%          for the registration form.
%
%  B. Stealthy Cooperative Jamming of UAV Networks: A Multi-Agent Generative
%     Approach
%       -> Same content, foregrounds stealth. Slightly punchier.
%
%  C. Breaking the Detection Assumption: Learned Cooperative Jamming against
%     Reactive UAV Anti-Jamming Systems
%       -> Names the actual contribution. Strongest for a paper title, but
%          commits to a result that has not been measured yet.
%
%  D. When Detection Fails: Limits of Reactive Anti-Jamming in UAV Networks
%       -> Defender-side framing. Useful if the emphasis shifts from attack
%          design toward evaluating existing countermeasures.
%
%  E. Effectiveness versus Detectability in Cooperative UAV Jamming
%       -> Broadest and most method-agnostic. Use if MARL might be dropped.
%==============================================================================

\documentclass[10pt,a4paper]{article}

\usepackage[utf8]{inputenc}
\usepackage[T1]{fontenc}
\usepackage{lmodern}
\usepackage[margin=2.1cm]{geometry}
\usepackage{amsmath,amssymb}
\usepackage{multicol}
% hyperref is optional: guarded so the file also compiles on minimal TeX installs.
\IfFileExists{etoolbox.sty}{\usepackage[hidelinks]{hyperref}}{}

\setlength{\parskip}{0.22em}
\setlength{\parindent}{0pt}

% compact section headings (no titlesec dependency)
\makeatletter
\renewcommand\section{\@startsection{section}{1}{\z@}%
  {0.9ex \@plus .2ex \@minus .2ex}{0.4ex \@plus .2ex}%
  {\normalfont\large\bfseries}}
\makeatother

\title{\vspace{-1.5em}\textbf{Cooperative Multi-Agent Generative Jamming\\
of UAV Networks under Detection Constraints}\\[0.3em]
\normalsize Bachelor's Thesis Project Proposal}

\author{\normalsize Rahul Rahman \\[-0.2em]
\small Supervisor: Dr.\ Antonio Di Maio, ETH Zurich}

\date{\small\today}

\begin{document}
\maketitle
\vspace{-2em}

%==============================================================================
\section{Introduction}
%==============================================================================

Unmanned aerial vehicle (UAV) networks are increasingly deployed as relays, sensor swarms and
mobile ad-hoc infrastructure, and jamming remains their dominant physical-layer
threat~\cite{pirayesh2022survey}. The countermeasure literature has converged on learned,
frequently multi-agent responses: repositioning relay UAVs away from an interferer, coordinated
spectrum access and frequency hopping, and joint trajectory and power
adaptation~\cite{nguyen2025marl,abolhassani2025coordinated}. These methods are effective, and
share a structural property that is rarely examined: they are \emph{reactive}. Each is triggered
by, or conditioned on, an observation that jamming is present --- a detector output, a classified
jammer type, or an observed collapse in link quality serving the same purpose.

That trigger is treated as reliable because the detection literature reports it as solved:
convolutional classifiers on spectrograms of the received signal achieve jamming detection
accuracies above 99\,\% in exactly the OFDM/UAV setting these defenses
assume~\cite{li2022jamming}. The resulting pipeline --- detect, classify, then relocate, hop or
re-power --- inherits a single point of failure. If an attacker degrades the link while remaining
statistically indistinguishable from the un-jammed condition, the countermeasure is never
invoked: the defense does not degrade gracefully, it simply never activates, and the network
experiences the degradation as ordinary channel impairment. The exceptions are \emph{proactive}
schemes that do not condition on detection at all --- uncoordinated frequency
hopping~\cite{strasser2009novel} and proactive data-driven
relocation~\cite{leuenberger2025proactive} being the clearest --- which makes them the most
informative comparison points, since they should retain performance where reactive schemes lose
it entirely.

Whether such an attack is constructible is an open question, and this thesis treats it as one.
Effectiveness and detectability are coupled by physics --- interference that flips bits
necessarily perturbs the received constellation --- but loosely enough to leave design margin,
since interference displacing symbols \emph{across} their decision boundaries rather than along
them causes the same errors at lower energy, and lower energy is harder to detect. Two choices
exploit that margin. \emph{Generative} policies parameterize the distribution of the transmitted
interference directly~\cite{goodfellow2014gan,durkan2019nsf} instead of selecting from a discrete
action set, which is what allows a signature to be shaped rather than a channel merely chosen ---
the mechanism underlying generative covert communication~\cite{wen2025generative}.
\emph{Multi-agent} coordination is used because splitting power and phase across several jammers
so their contributions add at the victim while each stays individually below threshold has no
closed-form solution, and because cooperative reinforcement learning already works for the
jamming task itself~\cite{valianti2024cooperative}. The result is an evasion attack against a
deployed classifier under constraints the adversarial machine learning literature does not
impose~\cite{sadeghi2019adversarial}: the perturbation must be physically realizable through the
attacker's own channel, and must raise the error rate rather than merely flip a label.

%==============================================================================
\section{Research Questions}
%==============================================================================

\textbf{RQ1.} Can a cooperative multi-agent, generative jamming policy degrade a UAV link while
remaining below the detection threshold of a state-of-the-art learned detector, and what
effectiveness--detectability trade-off does it achieve relative to baselines? The baselines span
the full range of attacker capability: barrage interference, a closed-form minimum-energy
attack, a single-agent learned attacker, and an omniscient cancellation attacker as an upper
bound.

\textbf{RQ2.} How much of the reactive anti-jamming stack still functions under such an attack?
This is measured at the level of the \emph{countermeasure} rather than the detector --- whether
the trigger fires, how long it takes, and what network-level performance results --- against
proactive, detection-free defenses as the comparison.

\textbf{RQ3.} What does adaptation cost either side? Treating the interaction as a round-based
offline game --- frozen detector, attacker optimized against it, detector retrained on those
attacks, attacker re-optimized --- this asks what re-closing the gap costs the defender in
accuracy and false-alarm rate, and how much of the loss the attacker subsequently recovers.

%==============================================================================
\section{Methodology}
%==============================================================================

\textbf{Minimal model first.} The work starts from the smallest model in which these questions
remain meaningful and adds one layer at a time, each addition justified by a question the
previous layer could not answer --- not for simplicity's sake, but because in a small model the
defender's optimal decision rule is analytically available, which upgrades ``this trained
detector fails to see the attack'' into the far stronger ``no detector can do better''. The base
model is a single QPSK link over an additive white Gaussian noise channel attacked by one
interferer under a hard power budget, with detection at the receiver from the received signal
alone. Noise is swept starting from the noiseless case: without noise the un-jammed constellation
is a finite set of exact points and any perturbation is certainly detectable, so the stealth
region exists only as a consequence of noise. Extensions follow in order: several coordinating
jammers with distinct link gains and phases; imperfect attacker synchronization, which weakens
the attacker and strengthens any surviving result; a multicarrier waveform over a
frequency-selective fading channel; and finally UAV mobility and network-level evaluation.

\textbf{Attacker.} The attack policy is generative and trained under centralized training with
decentralized execution~\cite{lowe2017maddpg} through a multi-agent-native environment
interface~\cite{terry2021pettingzoo}. Its objective is the induced error rate penalized by
detection probability, with the transmit-power budget imposed as a hard environment constraint
rather than a reward term. The training/execution asymmetry is stated explicitly as a modelling
assumption: the error rate is available to a centralized critic during training, whereas a
deployed jammer observes only its own transmission, its own channel estimate, and at most an
acknowledgement signal from the victim.

\textbf{Defenders and baselines.} The defender is a suite rather than a single model: an energy
detector at fixed false-alarm rate, a learned spectrogram classifier following the current state
of the art~\cite{li2022jamming}, and --- where the model admits it --- the optimal
likelihood-ratio test as reference. Above the detector sits the countermeasure itself, so that
RQ2 can be answered: a reactive, detection-triggered response (relocation or spectrum
reallocation~\cite{nguyen2025marl,abolhassani2025coordinated}) and a proactive, detection-free
one~\cite{leuenberger2025proactive,strasser2009novel}, evaluated side by side under the same
attack.

\textbf{Evaluation.} Attack performance is reported as bit and symbol error rate against
detection probability at a fixed false-alarm rate, under two conventions: strategies are compared
at \emph{matched detectability} rather than at equal transmit power, since an attack that raises
the error rate by raising its own signature is not an improvement; and every figure carries the
full reference envelope from the no-attacker floor to the omniscient ceiling. Defense performance
is reported as trigger rate, time-to-trigger and network throughput. Parameter studies cover
noise level, power budget, the number of jammers and legitimate nodes, and the exploitable
structure in the transmitted sequence. Experiments use Sionna~\cite{hoydis2022sionna} for the
physical layer and PyTorch for the learned components.

\textbf{Risks.} The principal risk is that the learned attacker does not separate from the
closed-form minimum-energy attack on a single link --- plausible, since against an independent
and uniformly distributed payload that perturbation is available analytically and a learner can
at best rediscover it. This is anticipated rather than avoided: the corresponding ablation is
scheduled early, and a negative result redirects the work toward the coordinated multi-jammer
setting, where no closed form exists and which is in any case the more relevant threat model for
UAV swarms. A second risk --- results that look strong only because the attacker was granted
unrealistic knowledge --- is answered by the assumption ladder, in which capabilities are removed
in defined steps.

\textbf{Work plan.} \textbf{(1)}~System, defender and threat model specification.
\textbf{(2)}~Minimal model: link, attacker ladder, detector suite, optimal-test reference.
\textbf{(3)}~Effectiveness--detectability characterization against the baselines (RQ1).
\textbf{(4)}~Generative and multi-agent attacker; comparison at matched detectability (RQ1).
\textbf{(5)}~Countermeasure-level evaluation, reactive versus proactive (RQ2).
\textbf{(6)}~Adaptation-cost rounds and synchronization realism (RQ3). \textbf{(7)}~Multicarrier,
fading and mobility; writing. Phases 1--3 are prerequisites; 5--7 are ordered by expected value
and may be truncated from the end.

%==============================================================================
{\scriptsize
\begin{multicols}{2}
\begin{thebibliography}{99}
\setlength{\itemsep}{0pt}
\setlength{\parsep}{0pt}
%==============================================================================

\bibitem{pirayesh2022survey}
H.~Pirayesh and H.~Zeng, ``Jamming attacks and anti-jamming strategies in wireless networks: A
comprehensive survey,'' \emph{IEEE Communications Surveys \& Tutorials}, 2022.

\bibitem{nguyen2025marl}
T.~D. Nguyen \emph{et al.}, ``Multi-agent deep reinforcement learning for collaborative UAV relay
networks under jamming attacks,'' 2025.

\bibitem{abolhassani2025coordinated}
B.~Abolhassani \emph{et al.}, ``Coordinated anti-jamming resilience in swarm networks via
multi-agent reinforcement learning,'' \emph{arXiv:2512.16813}, 2025.

\bibitem{li2022jamming}
Y.~Li \emph{et al.}, ``Jamming detection in OFDM-based UAV networks via spectrogram-tailored
machine learning,'' \emph{IEEE Access}, 2022.

\bibitem{strasser2009novel}
M.~Strasser, \emph{Novel Techniques for Thwarting Communication Jamming in Wireless Networks},
2009.

\bibitem{leuenberger2025proactive}
S.~Leuenberger, A.~Di Maio, and T.~Braun, ``Proactive decentralized multi-agent data-driven node
relocation for jammed mobile networks,'' in \emph{Wireless Days}, 2025.

\bibitem{valianti2024cooperative}
P.~Valianti \emph{et al.}, ``Cooperative multi-agent jamming of multiple rogue drones using
reinforcement learning,'' \emph{IEEE Trans. Mobile Computing}, 2024.

\bibitem{goodfellow2014gan}
I.~J. Goodfellow \emph{et al.}, ``Generative adversarial nets,'' in \emph{NeurIPS}, 2014.

\bibitem{durkan2019nsf}
C.~Durkan \emph{et al.}, ``Neural spline flows,'' in \emph{NeurIPS}, 2019.

\bibitem{wen2025generative}
Y.~Wen \emph{et al.}, ``Generative adversarial network-aided covert communication for cooperative
jammers in CCRNs,'' \emph{IEEE Trans. Information Forensics and Security}, 2025.

\bibitem{sadeghi2019adversarial}
M.~Sadeghi and E.~G. Larsson, ``Adversarial attacks on deep-learning based radio signal
classification,'' \emph{IEEE Wireless Communications Letters}, 2019.

\bibitem{lowe2017maddpg}
R.~Lowe \emph{et al.}, ``Multi-agent actor-critic for mixed cooperative-competitive
environments,'' in \emph{NeurIPS}, 2017.

\bibitem{terry2021pettingzoo}
J.~K. Terry \emph{et al.}, ``PettingZoo: Gym for multi-agent reinforcement learning,'' in
\emph{NeurIPS}, 2021.

\bibitem{hoydis2022sionna}
J.~Hoydis \emph{et al.}, ``Sionna: An open-source library for next-generation physical layer
research,'' \emph{arXiv:2203.11854}, 2022.

\end{thebibliography}
\end{multicols}
}

\end{document}
